\documentclass[12pt,a4paper]{article}

\usepackage[utf8]{inputenc}
\usepackage[T1]{fontenc}
\usepackage[brazil]{babel}
\usepackage{lmodern}
\usepackage{geometry}
\usepackage{microtype}
\usepackage{array}
\usepackage{tabularx}
\usepackage{enumitem}
\usepackage{xcolor}
\usepackage{fancyhdr}

\geometry{top=1.7cm,bottom=1.7cm,left=1.8cm,right=1.8cm}
\setlength{\parindent}{0pt}
\setlength{\parskip}{0.45em}
\setlist[itemize]{leftmargin=*,itemsep=0.25em,topsep=0.2em}
\definecolor{azul}{HTML}{173B57}
\definecolor{cinza}{HTML}{F2F4F7}

\pagestyle{fancy}
\setlength{\headheight}{15pt}
\fancyhf{}
\lhead{Sequência didática de \textit{machine learning}}
\rhead{Aula 5: árvore de decisão}
\cfoot{\thepage}
\renewcommand{\headrulewidth}{0.3pt}

\newcommand{\linha}{\rule{\linewidth}{0.35pt}}
\newcommand{\bloco}[1]{%
  \vspace{0.25em}
  \noindent\colorbox{cinza}{\parbox{0.96\linewidth}{#1}}%
  \vspace{0.3em}
}

\begin{document}

\begin{center}
  {\LARGE\bfseries\color{azul} Aula 5: árvore de decisão}\par
  \vspace{0.25em}
  {\large Largura, altura, regras e caminhos}\par
  \vspace{0.35em}
  {\small Roteiro de aplicação da sequência didática}\par
\end{center}

\bloco{\textbf{Lugar desta aula na pesquisa.} Este é o quinto encontro da sequência de oito aulas. A atividade apresenta uma classificação produzida por regras e caminhos de decisão. As produções individuais da aula têm função formativa; o professor registra dificuldades recorrentes e condições da aplicação.}

\section*{1. Identificação}

\begin{tabularx}{\linewidth}{@{}>{\bfseries}p{4.1cm}X@{}}
Turma: & \rule{8.5cm}{0.35pt} \\
Data: & \rule{8.5cm}{0.35pt} \\
Duração: & 50 minutos \\
Organização: & Exploração individual e cartão de síntese individual. \\
Recursos: & Plataforma no módulo de árvore de decisão, computadores do laboratório ou Chromebooks, quadro e cartão de síntese. \\
\end{tabularx}

\section*{2. Objetivos da aula}

Ao final do encontro, espera-se que os estudantes interpretem uma decisão como resultado de regras, percorram um caminho em uma árvore e reconheçam que aumentar a profundidade cria mais divisões, mas não assegura uma resposta melhor em qualquer situação.

\textbf{Relação curricular:} EM13CO10, EM13CO11 e EM13CO12.

\textbf{Conceitos em foco:} característica, regra, divisão, ramo, nó, folha, caminho de decisão, profundidade e classificação.

\section*{3. Preparação antes da aula}

\begin{itemize}
  \item Abrir o módulo de árvore de decisão e confirmar o cenário de objetos descritos por largura e altura.
  \item Conferir se as perguntas, os ramos e a categoria final podem ser vistos pela turma.
  \item Confirmar que cada estudante terá acesso individual ao módulo, em computador do laboratório ou Chromebook.
  \item Imprimir um cartão de síntese por estudante, sem código de participante e sem recolhimento como instrumento de pesquisa.
\end{itemize}

\bloco{\textbf{Cuidado didático.} A árvore organiza regras para a situação apresentada. Mais perguntas e mais ramos não significam, por si só, que a decisão será mais adequada, pois a qualidade depende dos dados, das regras escolhidas e da finalidade da tarefa.}

\section*{4. Desenvolvimento da aula}

\renewcommand{\arraystretch}{1.22}
\begin{tabularx}{\linewidth}{|>{\centering\arraybackslash}p{1.65cm}|p{5.0cm}|X|}
\hline
\textbf{Tempo} & \textbf{Ação com a turma} & \textbf{Mediação docente} \\
\hline
0--7 min & Retomar que classificações podem ser feitas de maneiras diferentes e apresentar a estratégia de decidir por perguntas. & Relacionar com a aula anterior: o k-vizinhos usa proximidade; a árvore organiza regras sobre características. \\
\hline
7--13 min & Apresentar objetos descritos por largura e altura. Pedir que cada estudante proponha uma regra simples para separar duas categorias. & Registrar algumas regras no quadro e perguntar qual informação seria usada primeiro. Evitar apresentar a árvore como decisão automática sem escolha humana. \\
\hline
13--18 min & Solicitar previsão sobre o efeito de aumentar a profundidade da árvore. & Perguntar se novas perguntas podem criar caminhos mais específicos e se isso garante uma resposta melhor para todos os casos. \\
\hline
18--33 min & Explorar o módulo guiado: divisão inicial, regras, ramos, folhas e profundidade. Acompanhar o caminho de um objeto até a categoria. & Relacionar cada ramo a uma condição e cada folha a uma saída. Não introduzir cálculos de impureza como exigência da aula. \\
\hline
33--43 min & Analisar um caso específico e comparar duas árvores com profundidades diferentes. & Pedir que a turma nomeie as perguntas percorridas e explique por que o objeto chegou à categoria indicada. \\
\hline
43--50 min & Preencher o cartão de síntese individual e realizar fechamento oral. & Circular pela turma, solicitar justificativas e usar os cartões apenas para orientar a conversa, sem recolhê-los como dado individual. \\
\hline
\end{tabularx}

\section*{5. Perguntas para conduzir a conversa}

\begin{itemize}
  \item ``Qual característica aparece na primeira pergunta?''
  \item ``Qual caminho esse objeto percorre na árvore?''
  \item ``Que regra leva a cada ramo?''
  \item ``O que muda quando a árvore fica mais profunda?''
  \item ``Ter mais regras sempre produz uma decisão melhor?''
\end{itemize}

\section*{6. Evidência a observar}

Observar se os estudantes explicam a saída como consequência de regras e dados de entrada, distinguindo o caminho da árvore de uma simples escolha por cor ou posição na tela. No registro docente, anotar dificuldades para ler condições, justificativas sobre profundidade e condições de uso da plataforma.

\newpage

\begin{center}
  {\LARGE\bfseries\color{azul} Cartão de síntese individual}\par
  \vspace{0.2em}
  {\large Aula 5: regras e caminhos de decisão}\par
  \vspace{0.35em}
  {\small Material formativo para discussão no final da aula}\par
\end{center}

\bloco{\textbf{Orientação.} Preencha o cartão individualmente. Ele serve para organizar a conversa da turma e não será usado como instrumento de pesquisa.}

\textbf{Nome do estudante:} \rule{8.5cm}{0.4pt}

\section*{Situação}

Uma árvore classifica objetos fictícios usando duas características: largura e altura. A primeira pergunta é se a largura é maior que um valor definido. Depois, cada ramo apresenta uma nova pergunta sobre altura antes de indicar uma categoria.

\section*{Questão 1}

Escreva uma regra possível para a primeira pergunta da árvore e indique o que acontece em cada resposta.

\vspace{0.25em}
\linha\par\vspace{0.75em}
\linha\par\vspace{0.75em}
\linha

\section*{Questão 2}

Escolha um objeto do exemplo e descreva o caminho que ele percorre até chegar a uma categoria.

\vspace{0.25em}
\linha\par\vspace{0.75em}
\linha\par\vspace{0.75em}
\linha

\section*{Questão 3}

Se a árvore receber mais uma pergunta antes de decidir, o que pode mudar? Isso garante uma classificação melhor em todos os casos? Explique.

\vspace{0.25em}
\linha\par\vspace{0.75em}
\linha

\end{document}
